% Transcription — fonds Alexandre Grothendieck, shelfmark 129, batch 6 (pages 101-114).
% Established from the facsimile archives/batches/129/batch-06.pdf, prepared
% from a copy of 129.pdf fetched through the project's relay
% (https://grothendieck-relay.onrender.com/source/129.pdf; PDF page k =
% archivists' page 100+k, no cover sheet), per the skill transcribe-grothendieck.
% The folder is 114 pages in six batches, transcribed in parallel by six
% passes; this is the sixth and last (pages 101-114, fourteen pages).
%
% Pass: Opus 5.5 (claude-opus-5-5), 2026-09-26 — first pass, unchecked against the pages by a human.
%
% Dating: \dating{[1970]} is the folder's own date in src/content/catalogue.ts,
% copied verbatim; since the folder has a date of its own, the group rule
% (« Descentes », 126-133) does not apply. Nothing on these leaves carries a
% date; page 112 is headed « Conférence à Nice » in his hand (noted where it
% stands, not used for dating).
%
% Numbering. \page{} follows the brief's mapping (PDF page k = 100+k), which
% is what keeps the facsimile in step. BUT the pencilled archivists' numbers
% bottom-left of these leaves read 103 ... 116 (PDF page k carries 102+k),
% whereas batch 4's first leaf carries 61 = 60+k. The pencil numbering thus
% runs two ahead somewhere in pages 81-100 (batch 5, being written in
% parallel; not present when this pass finished). FOR THE HUMAN CHECK:
% reconcile with the inventory's count of 114 pages.
%
% Pages skipped, without description: 102, 104, 108, 110 (typed pages of an
% unrelated text, versos of his sheets, blue duplicator copies numbered « n°
% 380 » / « n° 382 » — analysis, algebras of functions and analytic
% functions; nothing of this folder's subject; one handwritten word inserted
% on 102, hand undetermined, part of that text's own correction), and 111
% (a blank cream sheet).
%
% Typescripts: none of this folder's seminar typescripts occurs in this
% batch, and the ribbon typescript « Rappels sur la théorie de Dieudonné »
% begun at pages 79-80 does not continue here (it may continue in batch 5).
% No authorship decision was therefore needed or taken on it in this batch.
%
% Pages transcribed (all in hand): 101, 103, 105, 106, 107, 109, 112, 113,
% 114. Hands: 112-114 are a fair outline headed « Conférence à Nice » in
% the first person (« j'ai une idée »), which is the strongest evidence in
% the folder that a run is his; 101-109 are fast working notes in the same
% shorthand (« Kif-Kif », « cqfd »-style abbreviations, « loc. », « fid.
% plat »), taken as his. The exception is the head of page 106: four lines
% in English, in pencil, in a rounder hand (« ? Proof of 1st conj. »,
% « Galois-devissage »); whose they are cannot be told from the leaf, and
% the note there says so.
%
% His pagination: none.
%
% What the batch is about. (1) Pages 101-109: working notes on « objets de
% Tate » — a commutative group scheme G with an endomorphism f, f^n = 0,
% whose kernels Ker f^j are flat and f^{n-j} : G -> Ker f^j faithfully flat;
% fibrewise criterion (p. 101), proof and corollaries (unique sub-object and
% quotient of Tate of given height, p. 103); a lemma on a morphism of short
% exact sequences of such objects (p. 105); scattered computations with
% Ker f^j = Im f^{n-j} and a filtration picture (p. 106); a list of
% equivalent conditions a)-f) and their duals, then examples with Z/p^2Z,
% alpha_p, F, V, FV = p (p. 107); conditions for an epimorphism G -> H with
% kernel N to be formally smooth, via Phi(G) -> Phi(H) (p. 109). The objects
% are, in substance, truncated Barsotti-Tate groups of the typescript of
% chapter III (batches 3-4). (2) Pages 112-114: the outline of a talk,
% « Conférence à Nice », on Barsotti-Tate groups and crystals: definition of
% BT groups, Serre-Tate, the formal group, divided powers, the crystalline
% site (nilpotent and quasi-nilpotent p.d., exp and log), crystals in locally
% free modules (examples over W(A_0) and on a lifting over W with integrable
% connection), then the programme I)-V): the Dieudonné functor BT(S)° ->
% crystals, its F and V, full faithfulness, the classical case over a
% perfect field, the filtration 0 -> omega -> D(G)_S -> ... -> 0 and
% deformations as liftings of the filtration, and the isogeny category over
% a complete ring.
%
% Hardest pages: 106 (written in several orientations, crowded), 109 (the
% connective prose of conditions (i)-(iv bis)) and 114 (interlinear
% insertions). \ill{}: 53. \uncertain{}: 39.

\documentclass[11pt,a4paper]{article}
% Le préambule commun à toutes les transcriptions du fonds Grothendieck.
%
% Il tient en une idée : **l'incertitude se compose, elle ne se résout pas.**
% Une transcription qui lisse un mot douteux a détruit la seule chose pour
% laquelle on la faisait — un lecteur doit pouvoir savoir, à chaque endroit,
% ce qui était sur le feuillet et ce qui a été ajouté. Les six macros qui
% suivent sont donc l'appareil critique, et non de la décoration.
%
% Les mêmes commandes sont reconnues par `scripts/render.mjs`, qui produit la
% vue de lecture du site. Une macro ajoutée ici doit l'être aussi là-bas,
% faute de quoi le rendu échouera bruyamment — ce qui est voulu : un rendu
% silencieusement faux est pire qu'un rendu absent.

\ProvidesPackage{grothendieck}[2026/08/08 Transcriptions du fonds Grothendieck]

\usepackage{amsmath,amssymb,amsthm}
\usepackage{mathrsfs}
\usepackage{stmaryrd}
% Les diagrammes commutatifs sont partout dans ce fonds : c'est la langue
% même de la géométrie algébrique de Grothendieck, et une transcription qui
% les rendrait en prose ne transcrirait rien.
\usepackage{tikz-cd}
% Français par défaut — c'est la langue des feuillets — mais l'anglais est
% chargé aussi : la traduction et le résumé sont des documents anglais, et la
% typographie française (espaces avant les deux-points) y serait une faute.
% Ces documents basculent par \selectlanguage{english} en tête de corps.
\usepackage[english,french]{babel}
\usepackage{fontspec}
\usepackage{geometry}
\geometry{a4paper,margin=25mm,marginparwidth=28mm}
\usepackage{marginnote}
\usepackage{xcolor}
% ulem, not soul. `soul`'s \st and \ul are built on a character-by-character
% scan that XeLaTeX's font machinery breaks: the document fails with an
% undefined control sequence rather than degrading. ulem does the same two jobs
% and is Unicode-engine safe. `normalem` stops it hijacking \emph, which the
% transcriptions use for Grothendieck's own emphasis.
\usepackage[normalem]{ulem}
\usepackage{hyperref}
\hypersetup{colorlinks=true,linkcolor=black,urlcolor=black,pdfborder={0 0 0}}

% --- Métadonnées du lot -----------------------------------------------------
% Reprises telles quelles par le script de rendu, qui les affiche en tête de
% la vue de lecture. Elles fixent ce que le fichier prétend couvrir : sans
% elles, une transcription flotte sans point d'attache dans un fonds de seize
% mille feuillets.

\newcommand{\folder}[1]{\gdef\grothFolder{#1}}
\newcommand{\batch}[1]{\gdef\grothBatch{#1}}
\newcommand{\leaves}[2]{\gdef\grothFirst{#1}\gdef\grothLast{#2}}
\newcommand{\foldertitle}[1]{\gdef\grothFoldertitle{#1}}
\gdef\grothFolder{?}\gdef\grothBatch{?}\gdef\grothFirst{?}\gdef\grothLast{?}\gdef\grothFoldertitle{}

% --- Le résumé --------------------------------------------------------------
%
% Il ouvre la lecture modernisée, sous le titre « Résumé », et il est composé
% en petit. Ce n'est pas un ornement : c'est le seul passage du document qui
% ne soit pas des mathématiques, et un lecteur doit pouvoir le reconnaître
% avant de l'avoir lu — puis le sauter s'il connaît déjà le sujet, ou s'y
% arrêter s'il ne le connaît pas. Le filet qui le suit marque la reprise du
% corps.

\newenvironment{resume}
  {\small\setlength{\parskip}{0.45em}}
  {\par\medskip\hrule\medskip}

% --- Mots-clés ---------------------------------------------------------------
%
% En anglais, en clôture du résumé de la lecture modernisée : le vocabulaire
% moderne sous lequel chercher ce que les feuillets construisent. Le manifeste
% du site (`npm run manifest`) les extrait pour étiqueter la cote entière —
% cette ligne est la source unique des tags, il n'y a pas de fichier à côté.
\newcommand{\keywords}[1]{\par\smallskip\noindent
  {\footnotesize\textsc{Keywords} --- \textit{#1}}\par}

% --- L'appareil critique ----------------------------------------------------

% Le numéro de feuillet, en marge. C'est l'ancre qui permet de régler un
% désaccord sur une formule en regardant *un* feuillet plutôt qu'une chemise
% de six cents. Le site s'en sert aussi pour tourner les pages du fac-similé
% quand on fait défiler la transcription.
\newcommand{\leaf}[1]{%
  \marginnote{\footnotesize\color{gray!70}\textsf{#1}}%
  \label{leaf:#1}\ignorespaces}

% Illisible : on ne devine pas. Un mot inventé qui se lit comme les autres est
% le pire résultat possible.
\newcommand{\ill}{\textcolor{red!60!black}{[\,\ldots\,]}}

% Lecture douteuse : le mot est donné, et signalé comme incertain.
\newcommand{\uncertain}[1]{\uline{#1}}

% Ajout de l'éditeur — une lettre manquante, un mot sous-entendu.
\newcommand{\add}[1]{\textcolor{blue!50!black}{[#1]}}

% Biffé par Grothendieck lui-même. Ce qu'il a rayé fait partie du document :
% une rature dit souvent où la pensée a changé de direction.
\newcommand{\struck}[1]{\sout{#1}}

% Note du transcripteur, sur l'état matériel du feuillet.
\newcommand{\note}[1]{\par\smallskip{\small\color{gray!60!black}\textit{[#1]}}\par\smallskip}

% Note marginale de Grothendieck, distincte du corps du texte.
\newcommand{\marginal}[1]{\par\smallskip\noindent\hspace{1em}%
  {\small\color{gray!60!black}#1}\par\smallskip}

% --- Titre ------------------------------------------------------------------

\AtBeginDocument{%
  \begin{center}
    {\large\bfseries Fonds Alexandre Grothendieck --- cote n\textsuperscript{o}~\grothFolder}\\[2pt]
    {\itshape\grothFoldertitle}\\[4pt]
    {\small feuillets \grothFirst--\grothLast\ (lot \grothBatch)}\\[2pt]
    {\footnotesize\color{gray!60!black}Transcription établie d'après le fac-similé numérisé
     par l'Université de Montpellier.\\
     Les lectures incertaines sont soulignées, les illisibles notées [\,\ldots\,],
     les ajouts entre crochets.}
  \end{center}
  \vspace{1em}}

\endinput


\folder{129}
\batch{6}
\pages{101}{114}
\dating{[1970]}
\watermark{Édition de démonstration}
\foldertitle{Notes séminaire Montréal (Delale, Labute, Hakim, Messing) : copies de tapuscrits annotés (s.d.), notes manuscrites (s.d.)}

\begin{document}

\section{Objets de Tate}

\page{101}
\note{notes manuscrites à l'encre, d'une écriture rapide ; le haut de la
page commence en cours de démonstration, la page précédente de ce
raisonnement n'est pas dans ce lot}
\ill{} que c'est un isom.

$u$ mono $\Rightarrow$ $\mathrm{gr}^0(u)$ mono, car
$\mathrm{gr}^0(u) \simeq (f^{n-1}A \to f^{n-1}B)$

$u$ épi $\Rightarrow$ $\mathrm{gr}^0(u)$ épi, car
$\mathrm{gr}^0(u) \simeq (A/fA \to B/fB)$

$\mathrm{gr}^0(u)$ mono \uncertain{ou} épi $\Rightarrow$ $\mathrm{gr}^0(u)$ est
un isom, car $A/fA$ et $B/fB$ \uncertain{de même rang}.\note{souligné par lui}

$\mathrm{gr}^0(u)$ un isom $\Rightarrow$ $\mathrm{gr}(u)$ un isom, car
$\mathrm{gr}^i(u) \simeq \lbrace \mathrm{gr}^0(u) \rbrace^{\ill{}}$ \ill{}

$\mathrm{gr}(u)$ un isom $\Rightarrow$ $u$ un isom \uncertain{bien connu}.
\note{au-dessus de cette ligne, l'indication $0 \leq i \leq n-1$}

\emph{Cor}\note{« Cor » est écrit au-dessus d'un mot biffé, lu
\uncertain{Scol}} Pour tout $i$, $0 \leq i \leq n$, $f^iA$ est le seul
sous-objet de Tate\note{un « B » est écrit en exposant après « Tate »} de
$A$, de hauteur $n-i$, et de rang $\rho$ : en effet, un tel $B$ est contenu
dans $f^iA = \mathrm{Ker}\, f^{n-i}$, et d'après ce qui précède, il
\uncertain{lui} est égal.

\emph{Cor dual} : $A/f^{n-i}A$ est \struck{\ill{}} le seul objet quotient
de $A$ qui est de Tate, de hauteur $n-i$, de rang $\rho$.



$G$ schéma en groupes commutatif, \uncertain{plat} \struck{loc.} de
présentation finie sur $S$. Endom.\ $f$, et entier $n$ tel que $f^n = 0$.
Pour $1 \leq j \leq n-1$ fixé,

Conditions équivalentes :
\begin{enumerate}
  \item[(i)] [$\mathrm{Ker}\, f^j$ est plat, et] \struck{\ill{}}
  $f^{n-j} : G \to \mathrm{Ker}\, f^j$ est fid.\ \struck{plat}.
  \item[(ii)] $\forall s \in S$, $(G_s, f)$ est un objet de T.A de hauteur
  $n$ (dans la catégorie des \struck{\ill{}} groupes \uncertain{algébriques}
  commutatifs sur $k$).
\end{enumerate}
\note{« T.A » pour « Tate » ; les crochets autour de « Ker $f^j$ est plat,
et » sont de lui}

\page{103}
\emph{Dém.} Car (i) $\Rightarrow$ (ii), car (i) implique que
$\mathrm{Ker}\, f_s^j = \mathrm{Im}\, f_s^{n-j}$, et on conclut par le lemme
préliminaire.

\struck{(ii)} (ii) $\Rightarrow$ (i), car soit $H = \mathrm{Ker}\, f^j$,
considérons $G \xrightarrow{g} H$ induit par $f^{n-j}$ (existe, car
$f^n = 0$ i.e.\ $f^j(f^{n-j}) = 0$). Par l'hyp.\ (ii), il est
\uncertain{fid.} plat fibre par fibre. Comme $G$ est plat de prés.\ finie
sur $S$ et $H$ de prés.\ finie sur $S$, on en conclut que $g$ est
\struck{plat} fid.\ plat \emph{et} $H$ est plat,

\emph{Cor} Ceci ne dépend pas du choix de $j$, et les $\mathrm{Ker}\, f^j$
sont plats, \struck{indiqués par le $\mathrm{Ker}\, f^{j}/\mathrm{Ker}\, f^{j}$
\ill{}} et \struck{\ill{}} pour $j \geq j'$,
$\mathrm{Ker}\, f^j / \mathrm{Ker}\, f^{j'}$ est représentable, comme
schéma en groupes de type fini, et $f$ dans $G$ faisant \struck{\ill{}}
\ill{} \struck{\ill{}} \uncertain{plat}, de la catégorie de Tate de hauteur
$n$.\note{« (loc. cit.) » est écrit sous « hauteur $n$ », souligné et
relié par un trait à la ligne suivante ; la ligne commencée par « $f$ dans
$G$ faisant » est surchargée}

\emph{Cor} $G$ fini sur $S$ : rang de Tate en un pt. Il y a un seul
sous-objet de Tate de hauteur donnée $i \leq n$, de $\equiv$ rang $\rho$.

Kif-Kif pour quotients.

\section{Suites exactes d'objets de Tate}

\page{105}
$G' \subset G \subset G''$\note{ainsi en tête de page ; le $G$ du milieu
est repassé}

\begin{tikzcd}
0 \arrow[r] & G'_2 \arrow[r, hook] \arrow[d, "\alpha'"] & G_2 \arrow[r] \arrow[d, "\alpha"] & G''_2 \arrow[r] \arrow[d, "\alpha''"] & 0 \\
0 \arrow[r] & G'_1 \arrow[r, hook] & G_1 \arrow[r] & G''_1 \arrow[r] & 0
\end{tikzcd}

\note{à droite du diagramme, biffé d'un trait : « $0 \to \mathrm{Ker}\,
G''_1 \to$ »}

$p : G_2 \to G_2$ se factorise à travers $G_1$, et s'annule sur $G_1$,
induit un isom
\[
G_2/G_1 \xrightarrow{\ \sim\ } G_1 .
\]

$p : G'_2 \to G'_2$ se factorise : trouver $G'_1$ (\struck{sur} ou
\uncertain{OK} fibre par fibre), et s'annule sur $G'_1$ (\uncertain{même
raison}) donc induit
\[
G'_2/G'_1 \xrightarrow{\ \sim\ } G'_1 .
\]
\note{« OK » d'une autre encre, brune}

\begin{tikzcd}
G'_2 \arrow[r, hook] \arrow[d, no head] & G_2 \arrow[d] \\
G'_1 \arrow[r, hook] & G_1
\end{tikzcd}

\[
\mathrm{Ker}\, \alpha'' \simeq \underbrace{(G'_2 \cap G_1)}_{\mathrm{Ker}\, p_{G'_2}}/G'_1
\simeq \mathrm{Ker}\,(G'_2/G'_1 \to G'_1)
\]

\[
G'_2 = \overline{G'_{2\eta}}, \qquad
G'_1 = \overline{G'_{1\eta}} = \overline{G'_{2\eta} \cap G_{1\eta}}
\subset \overline{G'_{2\eta}} \cap \overline{G_{1\eta}}
= G'_2 \cap G_1 = \mathrm{Ker}(p_{G'_2})
\]
\note{au-dessus du $G'_2$ de l'avant-dernier membre, une surcharge
illisible ; l'indice $\eta$ y est corrigé}

De $\overline{G'_{2\eta} \cap G_{1\eta}}$ part un trait vers un encadré :
\begin{quote}
\uncertain{égalité} ssi $G'_2 \cap G_1$ plat$/S$
\end{quote}

Les questions dans \uncertain{ce cas} se \uncertain{ramènent} à
\emph{$\mathrm{Ker}\, \alpha''$}, qui est \uncertain{nul} \struck{\ill{}}
\struck{\ill{}} \ill{}.

\begin{tikzcd}[column sep=small]
0 \arrow[r] & \mathrm{Ker}\, \alpha'' \arrow[r] & G'_2/G'_1 \arrow[r] & G_2/G_1 \arrow[r] & G''_2/\mathrm{Im}\, G''_1 \arrow[r] & 0
\end{tikzcd}

\note{sous $G'_2/G'_1$, un signe $=$ vertical et $\widetilde{G}'_1$ ; sous
$G_2/G_1$, un $\simeq$ vertical et $G_1$ ; une flèche de $G'_1$ monte vers
$G_1$, un trait relie $G'_1$ à $\widetilde{G}'_1$}

\page{106}
\note{page de brouillon, écrite dans plusieurs sens ; on donne les
formules dans l'ordre où elles se laissent lire. Les quatre premières
lignes, en anglais et au crayon, sont d'une écriture plus arrondie que le
reste ; on ne peut pas dire d'après le feuillet si elles sont de lui}

\emph{? Proof of 1st conj.} :

(1)\note{chiffre entouré} Given
$0 \to \Phi' \to \Phi(G) \to \Phi'' \to 0$ (Galois-devissage), then can
define $G'$, $G''$ so $\Phi' = \Phi(G')$, $\Phi'' = \Phi(G'')$.



\struck{$f^{j}x$} \quad $f^j x = 0$, \quad $x = fy$, \quad
$f^{j+1}y = 0$, \quad $y = fz$, \quad $f^{j+1}z = 0$ ;
\quad $x = f^{n-j+1}t$, \quad $f^{n-1}t = 0$.

\[
\mathrm{Ker}\, f^j = \mathrm{Im} \qquad
\mathrm{Ker}\, f = \mathrm{Im}\, f^{n-1}
\]
\[
\mathrm{Ker}\, f^j = \mathrm{Im}\, f^{n-j}
\quad\overset{?}{\Longleftarrow}\quad
\mathrm{Ker}\, f^{j+1} = \mathrm{Im}\, f^{n-j-1}
\]
\note{entre les deux lignes, « $\supset$ \ill{} » et « $\subset$ \ill{} »,
puis un « $\Downarrow$ ? »}

$f^{j+1}x = 0$ \quad $\left\lbrace
\begin{array}{l}
fx \in \mathrm{Ker}\, f^j \\
fx = f^{n-j}y
\end{array}\right.$

\[
(f^j)^{-1}(f^{n-j-1}A) = f^{n-j-1}A
\]
\note{le membre de gauche est surchargé}

Écrits verticalement, dans la marge gauche et au milieu :
$f^{j-1}x = 0$, $f^{\ill{}}x = 0$, $x = f^{\ill{}}y$ ;
$A/f^{j}A \xrightarrow[\sim]{f} fA$,
$fA/f^2A \xrightarrow{\sim} f^2A$ ;
$(f^j)^{-1}(f^{\ill{}}A) = fA$ ; $\mathrm{Ker}\, f^{\ill{}}$ ;
pour $n = 3$, $f^3 = 0$, $\mathrm{Ker}\, f = \mathrm{Im}\, f^2$.

Dans la marge droite, verticalement :
$\mathrm{Ker}\, f = \mathrm{Im}\, f^{n-1} \Rightarrow
\mathrm{Ker}\, f^{j} = \mathrm{Im}\, f^{n-j}$, avec
$f^{j}x = 0$, $x = f^{n-j}y$, et
« i.e.\ $f^{n-j}\left[ f^{\ill{}} \ill{} \right] = 0$ » ;
\struck{\ill{}}\note{une ligne biffée à l'encre brune}.

Deux schémas de filtration, dessinés comme des segments gradués :
$A \supset fA \supset f^2A \supset \cdots \supset f^{n-j}A \supset
f^{n-j-1}A \supset \cdots \supset f^{n-1}A \supset f^nA = 0$,
et $f^{j}A \supset f^{j-1}A \supset \cdots \supset f^{n-1}A$, chacun avec
une accolade regroupant une partie des crans.
\note{dessins redécrits, non reproduits}

\page{107}
\emph{Conditions équivalentes}

\begin{itemize}
  \item[a)] $G'_1 = G'_2 \cap G_1$ \quad $(= \mathrm{Ker}\, p_{G'_2})$
  \item[b)] $G'_2 \cap G_1$ est plat$/S$
  \item[c)] $G'_2$ est un gpe de Tate de hauteur 2
  \item[d)] $G'_2/G'_1 \xrightarrow{\text{induit par } p} G'_1$ est
  \struck{\ill{}} \struck{\ill{}} un \emph{mono}\note{trois mots soulignés
  l'un sous l'autre, les deux derniers biffés de hachures ; on lit
  « mono » en tête, « épim » en dernier}
  \item[e)] $G'_2/G'_1 \xrightarrow{\text{induit par } p} G'_1$ est un
  \emph{épim.}
  \item[f)] $G'_2 \xrightarrow{p} G'_1$ est un \emph{épim}
  \item[f*)] $G''_1 \to G''_2$ est un mono.
  \item[e*)] \struck{\ill{}} $\mathrm{Ker}\,(G''_2 \xrightarrow{p} G''_1)
  \leftarrow G''_1$ est un mono, \uncertain{id} est un épim.
  \item[d*)]
  \item[c*)] $G''_1$ est un gpe de Tate de hauteur 2
  \item[b*)] …
  \item[a*)] …
\end{itemize}
\note{« f*) » corrige une lettre surchargée ; sous la flèche de e*), un mot
lu \uncertain{épim}}

\note{sous un trait horizontal, un encadré à gauche contient les notes
suivantes, dispersées}

\[
(\mathbb{Z}/p^2\mathbb{Z} \times \mathbb{Z}/p\mathbb{Z})_\eta \simeq G_{2\eta}
\]
$G_{\eta s}$ \uncertain{pureme} infinitésimal et unipotent.

$\mathbb{Z}/p^2\mathbb{Z}$, \quad $\mathbb{Z}/p^2\mathbb{Z} \to \alpha_p$,
\quad $G'_\eta \xrightarrow{p} G'_\eta$, \quad
$G'_s \xrightarrow{p} G'_s$, \quad $VF = p$.

$C \xrightarrow{u} C$, $u_s = F$ ; \quad
$C \xrightarrow{F} C^{(p)} \xrightarrow{V} C$ ; \quad
$C = C^{(p)}$ ? \quad $F$ \quad $FV = p$.

\[
V = uF, \qquad F = u^{-1}V, \qquad V^2 = up, \qquad F^2 = u^{-1}p
\]
\note{sous $F = u^{-1}V$, une ligne biffée : \struck{$VF = u$} ; un
triangle $C \xrightarrow{F} C$, $u \uparrow$, $C \xrightarrow{V}$ ; en
haut à gauche de l'encadré, un faisceau de droites concourantes}

\section{Morphismes formellement lisses}

\page{109}
\[
0 \to N \to G \xrightarrow{u} H \to 0
\quad \not\Longrightarrow \quad
\mathrm{Ker}\ \text{\uncertain{est}}\ \text{\uncertain{lisse}}
\]
\note{au-dessus de la fin de ligne, entouré : « formellement »}

\struck{$G(\hat{S}) \to H(\hat{S})$}

$G \to H$ \qquad $\Phi(G) \to \Phi(H)$ \emph{surj.}

$G \xrightarrow{u} H$ homom.\ de gpes \uncertain{invariable} \ill{} sur le
T.E.

Conditions équivalentes
\begin{itemize}
  \item[(i)] L'homom.\ $u(\ill{}) : G(\ill{}) \to H(\ill{})$ est un épim.
  [il suffit de \uncertain{vérifier} \ill{}]
  \item[(ii)] L'homom.\ $u(\ill{}) : G(\ill{}) \to H(\ill{})$ est un
  [\ill{}] épim \ill{}.
  \item[(iii)] L'homom.\ $u$ est formellement lisse et un épim.
  \item[(iii bis)] L'homom.\ $u_s : G_s \to H_s$ est formellt lisse et un
  épim.
  \item[(iv)] L'homom.\ de gpes \uncertain{ensemble} sur $k$ :
  $G_k \to H_k$ est \struck{\ill{}} [formellt lisse et] un épim et
  \[
  \Phi(G) \to \Phi(H) \quad \text{est surjectif}
  \]
  \item[(iv bis)]
\end{itemize}



\begin{itemize}
  \item[a)] Si $\Phi(u)$ est un isom, \ill{} $u$ est un isom
  ($\Phi$ est « conservatif »)
  \item[b)] Si $u$ est déjà formellt fid.\ plat, \struck{\ill{}} \ill{}
  [i.e.\ $\mathrm{Ker}\, u$ \uncertain{pas} formellt lisse] mais pas
  formellt lisse, alors $\Phi(u)$ pas surjectif…
\end{itemize}
\note{une accolade réunit a) et b) à gauche}

\section{Conférence à Nice}

\page{112}
\note{en haut à droite, encadré, de sa main : « Conférence à Nice » ;
titre souligné}

\emph{Groupes de Barsotti-Tate} [$p$ fixé]

1) $\mathrm{fppf}$ sur $S$.

$G$ groupe de B.T. si
\begin{itemize}
  \item[a)] $pG = G$
  \item[b)] $G$ de $p$-torsion
  \item[c)] ${}_pG = \mathrm{Ker}\, p\,\mathrm{id}_G$ fini loc.\ libre.
\end{itemize}

Alors les ${}_{p^n}G = G(n)$ le sont aussi : si $G(1)$ de rang $p^d$, alors
$G(n)$ de rang $p^{nd}$ ; on a $G = \varinjlim_n G(n)$.

Exemple important : $A$ schéma abélien,
\[
G(n) = {}_{p^n}A, \qquad G = {}_{p^\infty}A .
\]

\struck{2)} Th.\ de Serre-Tate
\[
S_0 \hookrightarrow S \qquad A \longmapsto (A_0, {}_{p^\infty}A, \varphi)
\]
une équivalence de catégories.\note{sous $S_0 \hookrightarrow S$, en petit,
lu \uncertain{défini de} \uncertain{com. $p$}}

2) Groupe formel $\widehat{G}$ associé à un groupe de B.T.

3) \emph{Cristaux}. Puissances divisées sur un idéal $I$ :
$x^{(n)}$ ($n \geq 0$) $\left(\frac{x^n}{n!}\right)$, \quad
$x^{(0)} = 1$, $x^{(1)} = x$, $x^{(n)} \in I$,
\[
\left\lbrace
\begin{array}{l}
(x+y)^{(n)} = \sum_{i+j=n} x^{(i)} y^{(j)} \\
(\lambda x)^{(n)} = \lambda^n x^{(n)} \\
\dots
\end{array}\right.
\]
$I^{(n)}$ idéal \struck{\ill{}} engendré par les monômes de poids $\geq n$
en les $x^{(i)}$ ($x \in I$).

\emph{Site cristallin absolu} \struck{\ill{}}

puiss.\ div.\ nilpotentes : $I^{(n)} = 0$ si $n$ grand

—— quasi-nilpotentes : $I^{(n)}$ de torsion pour $n$ grand, i.e.\
$(n-1)!\, I^{(n)} = 0$.

exponentielle : $\exp(x) = \sum_{n \geq 0} x^{(n)}$, \quad
$\log(1+x) = \sum_{n \geq 1} (-1)^{n-1} (n-1)!\, x^{(n)}$ \quad ok si p.d.\
nilp. ; ok —— quasi-nilp.

\emph{Site cristallin} (absolu) [\uncertain{relatif}] : \struck{\ill{}} On
exige que les p.d.\ \uncertain{soient compatibles aux} habituelles…
\note{sous la ligne, deux mots lus \uncertain{nilpot.} et « Berthelot »,
avec \ill{}}

\marginal{début de Théorie cohomologique (cf.\ Illusie) développée par
Berthelot}\note{écrit en oblique dans la marge inférieure gauche ; les mots
« début de » sont lus avec doute}

\page{113}
Cristaux de Modules [loc.\ libres] = Modules loc.\ libres sur le site
cristallin.

\emph{Ex.\ 1} $S$ \struck{schéma} $= \mathrm{Spec}\, A_0$, $A_0$ anneau
parfait de car.\ $p > 0$.

$A = W(A_0)$, alors les cristaux sur $A_0$ sont les Modules loc.\ libres
sur $A$.

\emph{Ex.\ 2} $X_0$ lisse sur corps parfait $k$, $W = W(k)$,
$\mathcal{X}$ relevant $X_0$ en schéma \uncertain{formel} lisse sur $W$,
\struck{[et $F : \mathcal{X} \to \mathcal{X}$ relevant Frobenius \ill{}
$\sigma_W$-linéaire de $\mathcal{X}$]}. Un cristal est un module loc.\
libre $M$ sur $\mathcal{X}$, avec connexion à courbure nulle rel$/W$.

Ce sont des genres d'animaux \struck{étudiés} rencontrés \uncertain{naturellement}
par Dwork, Katz, … \struck{\ill{}}



4)

I) $S$ schéma avec $p$ loc.\ nilpotent. On définit « foncteur de
Dieudonné »
\[
\mathrm{BT}(S)^{\circ} \longrightarrow \mathrm{Cris\,Mod\,loc\,lib}(S),
\qquad G \longmapsto \mathbb{D}(G)
\]
compatible avec chgt.\ de base.

II) \struck{\ill{}} $S$ de car.\ $p$, alors $M = \mathbb{D}(G)$ est muni de
$F_M$, $V_M$ avec
\[
F_M V_M = p\,\mathrm{id}, \qquad V_M F_M = p\,\mathrm{id}
\]
i.e.\ $\mathbb{D}(G)$ est un cristal de Dieudonné.

Il se pourrait que le foncteur
\[
\mathrm{BT}(S)^{\circ} \longrightarrow \mathrm{Cris\,Dieud}(S)
\]
soit pleinement fidèle, et j'ai une idée de ce que pourrait être
l'image essentielle…

\page{114}
III) Dans le cas où $S = \mathrm{Spec}(k)$, $k$ corps parfait de car.\
$p > 0$, on retrouve le foncteur de Dieudonné habituel.

IV) (Définitions…) \quad $\mathbb{D}(G)_S$ est filtré par
\[
0 \to \omega_{G^{\ill{}}} \to \mathbb{D}(G)_S \to \check{\omega}_{G^*} \to 0
\]
\note{les exposants du premier $\omega$ et de son indice sont surchargés}

Soit $S_0 \hookrightarrow S$ \quad \uncertain{défini} par \ill{} idéal
muni de p.\ div.\ ; supposons [que $p$ soit loc.\ nilpotent sur $S$, et]
que les p.\ div.\ \uncertain{soient} \ill{} \ill{}.

Pour les groupes $G_0$ tels que les fibres soient
\uncertain{composées} toroïdales, \uncertain{ou} que les fibres de $G_0$
soient connexes, prolonger $G_0$ en $G_0$ \uncertain{relevant}, i.e.\
donner $G$ qui prolonge $G_0$, « la même chose » que de prolonger la
filtration de $\mathbb{D}(G_0)_S$.

V) Prenons les groupes de BT à isogénie près. Supposons $A$ \uncertain{noethérien}
et complet pour la top.\ $(p\,A)$-adique, $A_0 = A/pA$. Alors foncteur
\[
G \longmapsto (G_0,\ \text{filtration sur}\ \mathbb{D}(G_0)_A)
\]
de $\mathrm{BTisog}(S) \longrightarrow \ill{}$ est pleinement fidèle.

\emph{Cas d'un anneau} \struck{de val.\ discrète} local complet
(\uncertain{à} corps résiduel parfait de car.\ $p > 0$ ; anneau de
val.\ discrète) : groupe de B.T.\ est connu quand on connaît $G_0$ sur $k$
(i.e.\ module de Dieud.\ $M$ sur $k$) et une filtration sur
$(M \otimes_W A)_p$.
\note{la parenthèse sur l'anneau local complet est une insertion
interlinéaire, reliée par un trait au mot « anneau »}

\end{document}
